\section[Une institution juridique]{Une institution au cœur de l’écosystème de la recherche juridique}

Depuis 1979, la Bibliothèque Cujas, devenue officiellement une bibliothèque interuniversitaire, est régie par un modèle de co-tutelle. Aux termes de la convention signée cette même année, l'Université Panthéon-Sorbonne en assure la gestion administrative, tandis que l'Université Panthéon-Assas participe activement au pilotage et à la direction de cet établissement \footcite{lobetCujasTraditionModernite2002} qui dessert aujourd'hui près de 16 500 usagers usagers inscrits à la bibliothèque. Le public consultant ses collections se composent \enquote{d'étudiants de [deuxième année de licence] (L2) aux enseignants-chercheurs et professionnels du droit} \footcite{BibliothequeInteruniversitaireCujas}. Cela représente un vaste public qui place la bibliothèque au coeur de l'enseignement du droit.

La Bibliothèque Cujas s'intègre dans différents réseaux documentaires nationaux lui permettant d'acquérir son statut de pôle d'excellence. À ce titre, elle intervient comme un pôle associé historique de la \ac{BNF}. La Bibliothèque Cujas participe au pilotage du \enquote{programme national de numérisation concertée en sciences juridiques}\footcite{Arabesques_69_12_dossier_balleypdfEspaceTravail}. En s'associant à la \ac{BNF} sur ce vaste programme, la Bibliothèque Cujas ambitionne d'accroître significativement la visibilité en ligne de ses collections. Ce partenariat ancien et solide vise ainsi à \enquote{créer un réseau documentaire thématique, rassemblant des institutions de tout statut et de toute taille qui mènent des projets en matière de numérisation du patrimoine juridique}\footcite{ProgrammeNationalNumerisation}. 

Lionel Maurel, conservateur à la \ac{BNF}, se charge de la phase prospection du projet afin d'identifier tous les \enquote{gisements documentaires en droit et [les] bibliothèques numériques existantes}\footcite{Arabesques_69_12_dossier_balleypdfEspaceTravail}. Par la suite, cette cartographie se concrétise sous la forme de \enquote{deux appels à initiatives pour la numérisation de corpus imprimés, articulés autour des quatre sources traditionnelles du droit : sources législatives et réglementaires, jurisprudence, doctrine, sources du droit local}\footcite{ProgrammeNationalNumerisation}. Grâce au soutien de la \ac{BNF}, les collections numérisées pour ces appels à projet bénéficient d'une excellente visibilité au sein de Gallica, la bibliothèque numérique de la \ac{BNF}. 

À cette époque, face à un contexte budgétaire particulièrement contraint, la Bibliothèque Cujas choisit de rationaliser ses coûts en s'appuyant sur l'infrastructure technique  de la \ac{BNF} \footcite{2007RefonteProgramme}. En faisant numériser certains de ses corpus sur les chaînes de numérisation de la \ac{BNF}, l'établissement bénéficie d'un canal de diffusion: ses collections sont directement intégrées dans Gallica \footcite{Arabesques_69_12_dossier_balleypdfEspaceTravail}. Cette collaboration permet à Cujas de tirer profit de l'immense visibilité de la bibliothèque numérique nationale pour valoriser ses propres fonds patrimoniaux, tout en garantissant théoriquement l'identification claire et de leur provenance institutionnelle \footcite{ProgrammeNationalNumerisation}. Toutefois, dans la pratique, cette co-diffusion se heurte à des limites importantes comme la dilution de l'identité institutionnelle ou le blocage lié au moissonnage.

Sur le plan technique, cette mise en commun documentaire exige un signalement de qualité des notices bibliographiques \footcite{coudrinDispositifsSignalementUniversite2014}. Cette collaboration étroite avec la \ac{BNF} s'articule concrètement autour de deux modalités de production distinctes, dont les programmes éditoriaux sont systématiquement définis et arbitrés par le Conseil scientifique et technique (CST) droit. La première modalité consiste en une numérisation directe des documents de la Bibliothèque Cujas sur les chaînes de production de la \ac{BNF}, comme ce fut le cas pour le chantier des journaux officiels coloniaux. Dans cette configuration, les fichiers et leurs métadonnées étant directement intégrés au réservoir de la \ac{BNF}, aucun protocole de moissonnage n'est requis. La seconde modalité repose sur un système de cofinancement : les deux institutions s'accordent sur un programme de numérisation, et la Bibliothèque Cujas réalise les opérations en interne grâce aux subventions allouées par la \ac{BNF}. C'est précisément dans ce second cas de figure que l'interconnexion technique entre les deux établissements repose sur le protocole de moissonnage \ac{OAI-PMH}, qui constitue une condition indispensable pour les bibliothèques numériques souhaitant rapprocher ces collections numériques dispersées \footcite{gauthierProtocoleOAIPMHFonctionnalitesa}. C'est par cette standardisation des échanges de données que la bibliothèque s'affranchit de l'isolement physique de ses magasins pour faire dialoguer ses collections juridiques avec les réservoirs scientifiques nationaux \footcite{stutzmannSignalerRessourcesNumerisees2013}.

La Bibliothèque Cujas s'intègre par ailleurs dans divers réseaux nationaux, au premier rang desquels figure le réseau CollEx-Persée. À ce titre, l'établissement est officiellement reconnu comme \enquote{bibliothèque délégataire CollEx en droit et sciences juridiques}, ce qui le positionne comme l'une des têtes du réseau national de l'Enseignement Supérieur et de la Recherche \ac{ESR}\footcite{BibliothequeInteruniversitaireCujas}. Afin de renforcer son interconnexion avec les autres bibliothèques, elle entreprend le signalement de ses fonds archivistiques, à l'instar du fonds ancien de l'ONU, au sein du catalogue national Calames. Enfin, Elle participe à l'alimentation de la base de données Le Doctrinal, où elle dépouille des revues scientifiques en droit dans le projet \ac{HOPPE-Droit} \footcite{HoppeDroit}. L'ensemble de ces initiatives renforce sa position centrale dans le paysage de la recherche juridique française.

La richesse de ses collections physiques renforce la volonté de la Bibliothèque Cujas de valoriser son patrimoine par la biais de la dématérialisation, afin de mieux s'adresser à son public de recherche. C'est dans cette perspective de préservation et de diffusion élargie qu'est entreprise la constitution d'un premier corpus de référence. Cette sélection méthodique d'ouvrages fondamentaux servira de socle technique et intellectuel à la création de sa première bibliothèque numérique, CujasNum, marquant ainsi son passage à l'ère du numérique.